% =============================================================================
% SECCION: ANALISIS DE COSTOS DE TRANSACCION
% =============================================================================

\subsection{Costos de Transacci\'on}
\label{sec:costs}

% ---------------------------------------------------------------------------
% 8.1  COSTOS EN ESTRATEGIAS ACTIVAS
% ---------------------------------------------------------------------------

\subsubsection{Modelo de Costos y Estructura}
\label{subsec:costs_intro}

Una cr\'itica recurrente a las estrategias de \textit{trading} activo es que los costos de transacci\'on erosionan la rentabilidad que muestran en simulaciones hist\'oricas \citep{barber2000}. El modelo de costos empleado en este estudio incorpora tres componentes (\hyperlink{glos:spread_bidask}{\textit{bid-ask spread}}, \hyperlink{glos:expense_ratio}{\textit{expense ratio}} y \hyperlink{glos:vol_drag}{\textit{volatility drag}}), calibrados de forma conservadora, con \textit{spreads bid-ask} que exceden los valores t\'ipicos de mercado, como se documenta m\'as adelante. El retorno neto diario se obtiene como

\begin{equation}
r_{\text{neto}} = r_{\text{bruto}} - c_{\text{expense}} - c_{\text{trading}} - c_{\text{drag}}
\label{eq:net_return}
\end{equation}

donde $r_{\text{bruto}} = r_f + p \cdot (r_{\text{mkt}} - r_f)$ con $p \in \{0, 1, 3\}$ seg\'un la posici\'on (CASH, SPY, UPRO). Los par\'ametros por instrumento se detallan en la Tabla~\ref{tab:tx_costs_detail}; UPRO incurre en un \textit{expense ratio} de 0.91\%/a\~no y un \textit{bid-ask spread} de 5~bps por transacci\'on, SPY en 0.09\%/a\~no y 2~bps, y la posici\'on en efectivo no aplica costo alguno. Conviene precisar que la exposici\'on apalancada se modela de forma sint\'etica, a saber, aplicando tres veces el retorno en exceso del SPY menos los costos descritos, en lugar de emplear la serie hist\'orica de precios del ETF UPRO; los \textit{benchmarks} pasivos (SPY y UPRO \textit{Buy \& Hold}) se costean con el mismo modelo, de modo que las comparaciones relativas no se ven favorecidas por un tratamiento asim\'etrico de costos.

El \textit{volatility drag} de los ETFs apalancados se modela seg\'un la aproximaci\'on de primer orden.

\begin{equation}
c_{\text{drag}} \approx \frac{1}{2}(L^2 - L)\,\sigma^2
\label{eq:vol_drag}
\end{equation}

\noindent donde $\sigma$ es la volatilidad diaria del subyacente. Para evitar subestimar el decaimiento en reg\'imenes turbulentos, $\sigma$ no se fija en un valor constante sino que se estima d\'ia a d\'ia mediante la volatilidad realizada de una ventana m\'ovil retrospectiva de 21 d\'ias, sin incorporar informaci\'on futura. Para UPRO ($L = 3$) con una volatilidad del orden de 1\% diario, el \textit{drag} ronda $0.03\%$/d\'ia ($\approx 7.5\%$/a\~no) y se amplifica en per\'iodos de alta volatilidad, constituyendo el componente de costo m\'as significativo para los modelos con exposici\'on apalancada frecuente.

% ---------------------------------------------------------------------------
% 8.2  DESGLOSE Y VALIDACION
% ---------------------------------------------------------------------------

\subsubsection{Desglose de Costos e Impacto}
\label{subsec:cost_breakdown}

La Tabla~\ref{tab:cost_breakdown} presenta el desglose completo para los 23 modelos. Los cuatro modelos ganadores exhiben ratios Tx/P\&L entre 15.8\% (RandomForest) y 39.0\% (CNN-LSTM), indicando que entre aproximadamente tres quintos y cinco sextos de su rendimiento bruto sobrevive a los costos. LSTM+Attention acumula \$13,459 en costos sobre un capital inicial de \$10,000, con un retorno neto del +395\% una vez descontadas esas fricciones. El costo promedio de los 23 modelos asciende a \$3,749 (37.5\% del capital inicial), lo que sugiere la conveniencia de modelar costos al evaluar estrategias activas.

Para verificar que los par\'ametros reflejan condiciones reales, se contrastan con datos reportados por los emisores \citep{proshares2026,statestreet2026} (Tabla~\ref{tab:cost_validation}). Los \textit{expense ratios} son pr\'acticamente id\'enticos a los vigentes (ratio $\approx 1.0\times$), pero los \textit{bid-ask spreads} empleados sobreestiman los costos reales de manera sustancial, dado que la mediana reportada por ProShares para UPRO es 2~bps (vs.\ 5~bps en este estudio, un factor de $\approx 2.5\times$) y la de SPY es $\approx$0~bps (vs.\ 2~bps). Los retornos netos reportados incorporan, por tanto, una penalizaci\'on por costos de ejecuci\'on significativamente mayor a la que enfrentar\'ia un operador real.

% ---------------------------------------------------------------------------
% 8.3  SENSIBILIDAD Y TURNOVER
% ---------------------------------------------------------------------------

\subsubsection{An\'alisis de Sensibilidad a Costos}
\label{subsec:cost_sensitivity}

La Tabla~\ref{tab:cost_sensitivity} muestra que los cuatro modelos l\'ideres mantienen rentabilidad positiva incluso con costos duplicados. LSTM+Attention registra +261\% con costos al doble, frente al +101\% del SPY, y su punto de \textit{break-even} se sit\'ua en 3.9$\times$ los costos actuales. RandomForest presenta el \textit{break-even} m\'as alejado (6.3$\times$, Tx/P\&L de 15.8\%), mientras que CNN-LSTM resulta m\'as sensible (2.6$\times$, Tx/P\&L de 39.0\%) aunque mantiene retorno positivo en todos los escenarios.

La Figura~\ref{fig:turnover_return} eval\'ua la hip\'otesis de que mayor \hyperlink{glos:turnover}{\textit{turnover}} se traduce en menores retornos netos; el coeficiente de determinaci\'on pr\'acticamente nulo ($R^2 = 0.00$) no la respalda, sugiriendo que la calidad de las se\~nales de entrada resulta m\'as relevante que la frecuencia de operaci\'on.

% ---------------------------------------------------------------------------
% 8.4  COSTOS NO CONSIDERADOS
% ---------------------------------------------------------------------------

\subsubsection{Limitaciones del Modelo de Costos}
\label{subsec:costs_not_modeled}

El modelo incorpora las fricciones consideradas m\'as relevantes para una estrategia \textit{Long-Only} basada en ETFs, pero una implementaci\'on real deber\'ia contemplar costos adicionales. El \textit{slippage} y el \textit{market impact} resultan insignificantes para capitales inferiores a varias decenas de millones de d\'olares, dado que SPY y UPRO figuran entre los ETFs con mayor volumen de negociaci\'on (\$35~mil millones y \$800~millones diarios, respectivamente); la restricci\'on de escalabilidad que ello implica para capitales mayores se retoma en la Secci\'on~\ref{sec:discussion}. Los impuestos sobre ganancias de capital representan otro costo no modelado que depende de la jurisdicci\'on fiscal del inversor pero afecta sim\'etricamente a todas las estrategias activas, por lo que las comparaciones relativas se mantienen v\'alidas. Los costos de infraestructura (servidores, datos, mantenimiento) constituyen costos fijos amortizables que no escalan con el tama\~no de las posiciones y no alteran la evaluaci\'on comparativa.
